% Pursuing Stacks (cotes 134-2 to 134-8): archive cote 134-4, site page 30 (PDF page 31), typescript page 118.
% typescript. Draft compared: ps3.tex:1477-1544.
% AI-transcribed by Claude Opus 5.5 (claude-opus-5-5, adaptive thinking, effort high), Message Batches API, batch ps2 (msgbatch_01NVwr1Yuo8qJf5rykeGBXPK), 2026-11-12; not yet checked by eye.
\scanpage{134-4/30}
\mspage{118}
and for general $C_h$ (satisfying the basic assumption $\unsure{C}_h \hand{\subset} \overline{C}_h$), $\overline{C}_h$ can be
viewed as the ``saturation'' of $C_h$ , namely the smallest saturated subset
of Ob(M) satisfying the basic assumption, or in other words, the
smallest \uline{contractibility} \uline{structure} on M\hand{,} \xout{\unsure{containing R}} such that the objects
in $C_h$ are contractible.

It may be worth while to state more explicitly the basic assumption
here\hand{,} and the saturation condition on $C_h$ .

(Bas 4) For any X in $C_h$ , we can find a finite sequqnee\xout{s} of maps
from X to X :
\[ f_o=id_X, f_1, \ldots, f_N=c_s \quad , \]
joining the identity map of X to a constant map $c_s$ (defined by some secti-%
on s of X ), in such a way that two consecutive maps $f_i,f_{i+1}$ are $C_h$-homo-%
topic\xout{x} in the strict sense\hand{\unsure{*}}, namely we can find $Y_i$ in $C_{\hand{h}}$ and two sections
$\hand{\delta}^i_o$ and $\hand{\delta}^i_1$ of $Y_i$ over e, and a map
\[ h_i : Y_i x X \longrightarrow X , \]
such that
\[ h_i\hand{\circ}(\hand{\delta}^i_o x id_X) = f_i \ , \quad h_i\hand{\circ}(\hand{\delta}^i_1 x id_X) = f_{i+1} \quad . \]
(Sat 4) \xout{For} \uline{Any} object X in M satisfying the condition just stated is
in $C_h$ .

The third dotted arrow can be viewed\hand{,} as \inserted{denoting,} an embedding of $Hom_4(M)$
(the set of all contractibility structures on M) into $Hom_3(M)$ , we finally
get a cascade of three inclusions
\[ Hom_4(M) \hookrightarrow Hom_3(M) \hookrightarrow Hom_2(M) \hookrightarrow Hom_1(M) \quad , \]
in terms of which a contractibility structure \inserted{on M} can be viewed as a particu-%
lar case of \inserted{any of} the three types of homotopy structures on M considered before.

If we start with a homotop\unsure{y}\xout{sm\illegible{1}} \inserted{structure} \inserted{$W_h$} on M , and consider the corres-%
ponding set $C_h$ of contractible \xout{elements} objects of M (namely objects X
such that $X \to e$ is in $W_h$), it is pretty clear that $C_h$ satisfies the satu-%
ration condition (Sat 4), but by no means clear that it satisfies the
basic assumption (Bas 4), even in the special case when we assume more\xout{ax}
over that $W_h$ comes from a weak homotopy interval structure $\hand{\Sigma}_h$ on M .
The trouble comes from the circumstance that there is no reason in general
that the contractibility of an object X of M can be described in terms of
\xout{homotopies} a sequence of elementary homotopies between maps $f_i : M \to M$
(joining $id_M$ to a constant map ) involving weak homotopy intervals $\underline{I}_i$
\uline{which are moeover contractible}. \xout{There is} I doubt this is always so, and
there does'nt come either any plausible extra condition on $\hand{\Sigma}_h$ which may
insure this, except \inserted{precisely} that $\hand{\Sigma}_h$ can be generated (through saturation)

\margin{30}
